\documentclass[12pt,a4paper]{article}
\usepackage[utf8]{inputenc}
\usepackage[spanish]{babel}
\usepackage{geometry}
\usepackage{array}
\usepackage{tabularx}
\usepackage{xcolor}
\usepackage{enumitem}
\usepackage{hyperref}
\usepackage{graphicx}
\usepackage{fancyhdr}
\usepackage{titlesec}
\usepackage{float}

\geometry{margin=2.5cm}
\pagestyle{fancy}
\fancyhf{}
\fancyhead[L]{\small Proyecto LabCircuitos}
\fancyhead[R]{\small Plan de Trabajo}
\fancyfoot[C]{\thepage}
\renewcommand{\headrulewidth}{0.4pt}

\definecolor{primary}{RGB}{59,130,246}
\definecolor{success}{RGB}{34,197,94}
\definecolor{warning}{RGB}{234,179,8}
\definecolor{error}{RGB}{239,68,68}
\definecolor{surface}{RGB}{26,27,46}
\definecolor{textlight}{RGB}{226,232,240}

\titleformat{\section}{\Large\bfseries\color{primary}}{}{0em}{}[\vspace{-0.3em}\rule{\textwidth}{0.5pt}]
\titleformat{\subsection}{\large\bfseries\color{primary!80!black}}{}{0em}{}

\begin{document}

\begin{titlepage}
    \centering
    \vspace*{3cm}
    {\Huge\bfseries\color{primary} LabCircuitos\\[0.5em]}
    {\Large Simulador de Circuitos Eléctricos Interactivo\\[2cm]}
    {\large\textbf{Plan de Trabajo}}\\[1cm]
    {\large 3 Integrantes · 4 Días · Corrección y Puesta a Punto}\\[2cm]
    \vfill
    {\large Fecha: \today}
\end{titlepage}

\section{Resumen del Proyecto}

\textbf{LabCircuitos} es un simulador de circuitos eléctricos en tiempo real con arquitectura frontend/backend. Utiliza React + TypeScript + Vite en el frontend y FastAPI + PySpice/Ngspice en el backend. El proyecto se encuentra en un estado funcional avanzado pero requiere corrección de errores, pulido de interfaz, y validación integral para alcanzar un nivel profesional.

\subsection*{Tecnologías Involucradas}
\begin{minipage}[t]{0.48\textwidth}
    \textbf{Frontend}\\
    React, TypeScript, Vite, pnpm, Tailwind CSS, React Flow, Zustand, Plotly, Framer Motion, Lucide React
\end{minipage}
\hfill
\begin{minipage}[t]{0.48\textwidth}
    \textbf{Backend}\\
    Python, FastAPI, PySpice, Ngspice, NumPy, SciPy, Uvicorn
\end{minipage}

\section{Objetivos}

\begin{enumerate}[label=\arabic*.]
    \item Corregir todos los bugs conocidos en frontend y backend.
    \item Asegurar que la simulación SPICE funcione correctamente con circuitos de prueba.
    \item Validar que componentes se arrastren, conecten, y respondan en tiempo real.
    \item Pulir la interfaz de usuario al nivel de un laboratorio profesional.
    \item Traducir y unificar toda la interfaz en español.
    \item Generar un build de producción sin errores.
\end{enumerate}

\section{Distribución de Tareas}

\subsection*{Roles}
\begin{itemize}
    \item \textbf{Luisa (Integrante A)} — Backend + Integración (SPICE, simulación, API, validación de circuitos)
    \item \textbf{Miguel (Integrante B)} — Frontend núcleo (editor, componentes, conexiones, store, hooks)
    \item \textbf{Josue (Integrante C)} — Frontend UI/UX + Supervisor General (toolbar, paneles, calculadora, gráficas, dark mode, revisión de calidad)
\end{itemize}

\subsection*{Cronograma}

\begin{table}[H]
\centering
\small
\begin{tabularx}{\textwidth}{|c|X|X|X|}
\hline
\textbf{Día} & \textbf{A (Backend)} & \textbf{B (Editor)} & \textbf{C (UI/UX)} \\
\hline
\textbf{Día 1} &
\textbullet~Verificar ngspice instalado\\&
\textbullet~Probar endpoint /api/simulate\\&
\textbullet~Corregir ground validator BFS\\&
\textbullet~Agregar tipos: diodo, transistor & 
\textbullet~Probar drag de componentes\\&
\textbullet~Corregir onNodesChange\\&
\textbullet~Verificar handles y conexiones\\&
\textbullet~Probar batería terminales +/− & 
\textbullet~Verificar dark mode global\\&
\textbullet~Probar sidebar y buscador\\&
\textbullet~Verificar traducción española\\&
\textbullet~Reportar bugs encontrados \\
\hline
\textbf{Día 2} &
\textbullet~Simulación automática al conectar\\&
\textbullet~Probar circuito 9V + R + LED\\&
\textbullet~Validar nodeVoltages arrays\\&
\textbullet~Corregir errores branchCurrents &
\textbullet~Corregir selección múltiple\\&
\textbullet~Probar undo/redo\\&
\textbullet~Verificar rotación y duplicado\\&
\textbullet~Probar multi-selección + delete &
\textbullet~Rediseñar SVGs restantes\\&
\textbullet~Probar calculadora (7 tabs)\\&
\textbullet~Verificar gráficas Plotly\\&
\textbullet~Probar panel de propiedades \\
\hline
\textbf{Día 3} &
\textbullet~Pruebas de estrés: 10+ componentes\\&
\textbullet~Medir tiempos de simulación\\&
\textbullet~Cache de resultados\\&
\textbullet~Manejo de errores robusto &
\textbullet~Probar LED on/off con simulación\\&
\textbullet~Verificar animación de cables\\&
\textbullet~Probar todas las herramientas\\&
\textbullet~Probar circuito demo completo &
\textbullet~Pulir layout y espaciado\\&
\textbullet~Agregar micro-animaciones\\&
\textbullet~Probar calculadora en oscuro\\&
\textbullet~Verificar responsividad \\
\hline
\textbf{Día 4} &
\textbullet~Pruebas de regresión backend\\&
\textbullet~Documentación API\\&
\textbullet~Últimas correcciones &
\textbullet~Build final sin errores\\&
\textbullet~Prueba integración completa\\&
\textbullet~Correcciones menores &
\textbullet~Revisión final UI\\&
\textbullet~Preparar presentación\\&
\textbullet~Generar reporte \\
\hline
\end{tabularx}
\end{table}

\section{Checklist de Validación}

Cada integrante debe marcar como listo cada ítem antes del Día 4.

\subsection*{Backend — Integrante A}
\begin{enumerate}[label=\textbf{A\arabic*}]
    \item \textbf{Simulación funcional} — POST /api/simulate retorna nodeVoltages y branchCurrents correctos.
    \item \textbf{Ground validator} — Detecta tierra conectada mediante BFS sobre el grafo del circuito.
    \item \textbf{Circuitos de prueba} — 9V + R + LED, RLC, divisor de voltaje.
    \item \textbf{Manejo de errores} — Circuito inválido retorna error descriptivo, no crash.
    \item \textbf{Rendimiento} — Simulación de 10+ componentes en <2 segundos.
    \item \textbf{Netlist correcta} — Generación SPICE bien formada.
\end{enumerate}

\subsection*{Editor — Integrante B}
\begin{enumerate}[label=\textbf{B\arabic*}]
    \item \textbf{Drag de componentes} — Todos los componentes se mueven libremente en el canvas.
    \item \textbf{Conexiones} — Cables se crean desde handles reales, siguen curvas suaves.
    \item \textbf{Batería} — Terminal + (rojo) y − (verde) funcionales y visibles.
    \item \textbf{LED} — Se enciende solo con simulación activa y corriente > 1e-6 A.
    \item \textbf{Selección} — Click selecciona, multi-selección con Shift, Delete elimina.
    \item \textbf{Undo/Redo} — Ctrl+Z deshace, Ctrl+Shift+Z rehace.
    \item \textbf{Rotación/Duplicado} — Botones funcionales en panel de propiedades.
    \item \textbf{Herramientas} — Select, Wire, Probe funcionan correctamente.
\end{enumerate}

\subsection*{UI/UX — Integrante C}
\begin{enumerate}[label=\textbf{C\arabic*}]
    \item \textbf{Dark mode} — Tema oscuro consistente en toda la aplicación.
    \item \textbf{Sidebar} — Categorías, buscador, componentes compactos con drag.
    \item \textbf{Barra superior} — Botones (Nuevo, Abrir, Guardar, Exportar), estado de simulación.
    \item \textbf{Panel derecho} — Multímetro y Propiedades con mediciones en tiempo real.
    \item \textbf{Gráficas} — Plotly con múltiples sondas, zoom, toggle, limpiar.
    \item \textbf{Calculadora} — 7 tabs funcionales con modo oscuro.
    \item \textbf{Traducción} — 100\% en español, sin términos en inglés.
    \item \textbf{SVGs} — Todos los componentes con diseño profesional y terminales.
\end{enumerate}

\section{Pruebas de Integración Final}

\subsection{Circuito de Prueba Completo}
\begin{enumerate}
    \item Cargar circuito demo desde la pantalla de bienvenida.
    \item Verificar que aparecen 6 componentes en el canvas.
    \item Arrastrar cada componente para confirmar movimiento libre.
    \item Hover sobre handles → deben aparecer puntos azules.
    \item Iniciar simulación → LED debe encenderse (amarillo brillante).
    \item Verificar que el voltaje del LED se muestra en el nodo.
    \item Pausar simulación → LED debe apagarse inmediatamente.
    \item Reanudar → LED debe encenderse de nuevo.
    \item Abrir calculadora → probar Ley de Ohm (9V, 0.011A, 818Ω).
    \item Verificar gráficas con datos reales del osciloscopio.
\end{enumerate}

\subsection{Criterios de Aceptación}
\begin{itemize}
    \item El build de producción se completa sin errores (\texttt{pnpm build}).
    \item El backend responde en <2 segundos para circuitos de até 15 componentes.
    \item La interfaz no muestra errores de consola en Chrome/Firefox.
    \item Todos los textos están en español.
    \item El LED responde correctamente al estado de simulación.
    \item Los componentes se arrastran sin bloqueos.
\end{itemize}

\section{Entregables}

\begin{enumerate}
    \item Código fuente completo (frontend + backend) en repositorio Git.
    \item Build de producción en \texttt{dist/}.
    \item Backend ejecutable con \texttt{python main.py}.
    \item Documentación de API (endpoints, request/response).
    \item Video de demostración (opcional).
    \item Este plan de trabajo cumplimentado.
\end{enumerate}

\vfill
\begin{center}
\rule{0.6\textwidth}{0.5pt}\\[0.5em]
{\small Documento generado el \today}
\end{center}

\end{document}
